\section{梯度下降与高级优化器}

\subsection{梯度下降算法}

梯度下降是训练所有深度学习模型的基础算法：

$$W \leftarrow W - \alpha \cdot \nabla_W L$$

\begin{itemize}
    \item $\alpha$：学习率（步长），控制每一步走多远
    \item $\nabla_W L$：损失函数对权重的梯度
    \item 负号：沿梯度反方向走（下坡）
\end{itemize}

\begin{figure}[H]
\centering
\includegraphics[width=0.95\textwidth]{slides-images/slide-027.jpg}
\caption{梯度下降可视化：从初始点出发，沿负梯度方向逐步逼近最小值\protect\footnotemark}
\end{figure}
\footnotetext{Slides 第31页。}

\begin{knowledgebox}{何时停止训练？}
梯度下降的迭代过程何时结束？通常有两种策略：（1）设定固定的迭代次数；（2）设定停止准则，当损失的减少量低于某个阈值（如 $10^{-5}$）时停止。实践中更常用固定迭代次数。
\end{knowledgebox}

\subsection{随机梯度下降（SGD）}

标准梯度下降需要在整个训练集上计算损失和梯度，当数据量很大时计算代价极高。\textbf{随机梯度下降}（Stochastic Gradient Descent, SGD）通过在每步只使用一个小批量（mini-batch）来近似完整梯度：

\begin{lstlisting}[caption={SGD 的基本实现}]
for t in range(num_iterations):
    mini_batch = sample_data(X_train, batch_size=256)
    dW = compute_gradient(mini_batch, W)
    W = W - learning_rate * dW
\end{lstlisting}

\begin{figure}[H]
\centering
\includegraphics[width=0.95\textwidth]{slides-images/slide-029.jpg}
\caption{随机梯度下降：每次只使用数据子集计算梯度\protect\footnotemark}
\end{figure}
\footnotetext{Slides 第33页。}

\begin{importantbox}{Epoch 的概念}
实践中不是完全随机采样，而是将数据集打乱顺序后依次遍历。遍历完整个数据集一次称为一个 \textbf{epoch}。多个 epoch 的训练可以让模型反复学习所有数据。
\end{importantbox}

\subsection{SGD 的问题}

\begin{figure}[H]
\centering
\includegraphics[width=0.95\textwidth]{slides-images/slide-031.jpg}
\caption{SGD 在窄谷（高条件数）中的震荡问题：垂直方向梯度大导致来回震荡，水平方向进展缓慢\protect\footnotemark}
\end{figure}
\footnotetext{Slides 第35页。}

SGD 存在四个主要问题：

\begin{enumerate}
    \item \textbf{高条件数/窄谷}：梯度在陡峭方向很大、平坦方向很小，导致来回震荡而非直接走向最优点
    \item \textbf{局部极小值}：可能陷入局部最优，无法继续下降
    \item \textbf{鞍点}：在高维空间中更常见——梯度为零但不是极值点
    \item \textbf{梯度噪声}：因为只用了数据子集，每步梯度是有噪声的
\end{enumerate}

\begin{figure}[H]
\centering
\includegraphics[width=0.95\textwidth]{slides-images/slide-032.jpg}
\caption{局部最小值和鞍点：SGD 可能在这些点停滞，因为梯度为零或接近零\protect\footnotemark}
\end{figure}
\footnotetext{Slides 第36页。}

\begin{warningbox}{鞍点在高维空间中比局部极小值更常见}
研究表明，随着模型参数维度增加，鞍点出现的频率远高于局部极小值。在鞍点处，某些方向是``谷底''而另一些方向是``山顶''，梯度在所有方向上为零，导致优化停滞。这是大型神经网络训练中的一个严重实际问题。
\end{warningbox}

\subsection{SGD + Momentum}

动量方法借鉴了物理学中滚球的概念：球在下坡时会积累速度，即使遇到平坦区域也能凭借惯性继续前进。

\begin{figure}[H]
\centering
\includegraphics[width=0.95\textwidth]{slides-images/slide-035.jpg}
\caption{动量如何解决 SGD 的三个问题：越过局部极小值、穿越鞍点、减少窄谷中的震荡\protect\footnotemark}
\end{figure}
\footnotetext{Slides 第39页。}

数学形式如下：

$$\mathbf{v}_t = \rho \cdot \mathbf{v}_{t-1} + \nabla_W L$$
$$W = W - \alpha \cdot \mathbf{v}_t$$

\begin{itemize}
    \item $\mathbf{v}_t$：速度向量（过去梯度的加权平均）
    \item $\rho$：动量系数（通常取 0.9），控制历史梯度的权重
    \item 更新方向不再是当前梯度，而是速度向量
\end{itemize}

\begin{importantbox}{动量解决 SGD 三大问题}
\begin{enumerate}
    \item \textbf{局部极小值}：积累的速度使球能``滚出''浅的局部极小值
    \item \textbf{鞍点}：历史动量推动模型继续穿越平坦区域
    \item \textbf{窄谷震荡}：垂直方向的来回震荡在速度累加中相互抵消，而水平方向的一致运动不断累积
    \item \textbf{噪声}：多步梯度的平均效果减少了单步噪声的影响
\end{enumerate}
\end{importantbox}

\subsection{RMSProp}

RMSProp 由 Hinton 在 2012 年提出，核心思想是\textbf{逐元素缩放梯度}——在梯度大的方向走小步，在梯度小的方向走大步。

$$\mathbf{s}_t = \beta \cdot \mathbf{s}_{t-1} + (1-\beta) \cdot (\nabla_W L)^2$$
$$W = W - \frac{\alpha}{\sqrt{\mathbf{s}_t} + \epsilon} \cdot \nabla_W L$$

\begin{figure}[H]
\centering
\includegraphics[width=0.95\textwidth]{slides-images/slide-042.jpg}
\caption{RMSProp 的代码实现：通过梯度平方的运行平均来缩放学习率\protect\footnotemark}
\end{figure}
\footnotetext{Slides 第46页。}

\begin{itemize}
    \item $\mathbf{s}_t$：梯度平方的指数移动平均
    \item 陡峭方向：$\mathbf{s}_t$ 大，除以大数 $\rightarrow$ 步长变小
    \item 平坦方向：$\mathbf{s}_t$ 小，除以小数 $\rightarrow$ 步长变大
\end{itemize}

\begin{figure}[H]
\centering
\includegraphics[width=0.95\textwidth]{slides-images/slide-043.jpg}
\caption{SGD、SGD+Momentum、RMSProp 的优化轨迹对比：RMSProp 能快速转向平坦方向\protect\footnotemark}
\end{figure}
\footnotetext{Slides 第47页。}

\subsection{Adam 优化器}

\textbf{Adam}（Adaptive Moment Estimation）是现代深度学习中\textbf{最常用}的优化器，它结合了 Momentum 和 RMSProp 的优点：

\begin{figure}[H]
\centering
\includegraphics[width=0.95\textwidth]{slides-images/slide-044.jpg}
\caption{Adam 优化器：结合一阶矩（动量）和二阶矩（自适应学习率），并加入偏差修正\protect\footnotemark}
\end{figure}
\footnotetext{Slides 第48页。}

$$\mathbf{m}_t = \beta_1 \cdot \mathbf{m}_{t-1} + (1-\beta_1) \cdot \nabla_W L \qquad \text{（一阶矩/动量）}$$
$$\mathbf{v}_t = \beta_2 \cdot \mathbf{v}_{t-1} + (1-\beta_2) \cdot (\nabla_W L)^2 \qquad \text{（二阶矩/自适应缩放）}$$
$$W = W - \frac{\alpha \cdot \hat{\mathbf{m}}_t}{\sqrt{\hat{\mathbf{v}}_t} + \epsilon}$$

\begin{warningbox}{Adam 的初始步长问题}
Adam 的 $\beta_1, \beta_2$ 通常初始化为接近 1 的值（如 0.9, 0.999），而 $\mathbf{m}_0 = \mathbf{v}_0 = 0$。这意味着在第一步中，$\mathbf{v}_1$ 会非常接近零，导致除以极小数，产生一个不合理的大步长。为此 Adam 引入了\textbf{偏差修正}（bias correction），通过 $\hat{\mathbf{m}}_t = \mathbf{m}_t / (1 - \beta_1^t)$ 和 $\hat{\mathbf{v}}_t = \mathbf{v}_t / (1 - \beta_2^t)$ 来补偿初始时刻的估计偏差。
\end{warningbox}

常用默认参数：$\beta_1 = 0.9$，$\beta_2 = 0.999$，$\epsilon = 10^{-8}$，$\alpha = 10^{-3}$。

\begin{figure}[H]
\centering
\includegraphics[width=0.95\textwidth]{slides-images/slide-047.jpg}
\caption{四种优化器的收敛轨迹对比：Adam 结合了 Momentum 的加速性和 RMSProp 的方向修正能力\protect\footnotemark}
\end{figure}
\footnotetext{Slides 第51页。}

\subsection{AdamW：解耦权重衰减}

\begin{figure}[H]
\centering
\includegraphics[width=0.95\textwidth]{slides-images/slide-048.jpg}
\caption{Adam vs AdamW：AdamW 将 L2 正则化与动量计算分离\protect\footnotemark}
\end{figure}
\footnotetext{Slides 第52页。}

标准 Adam 在计算梯度时包含 L2 正则化项，这意味着正则化会影响动量和自适应学习率的计算。\textbf{AdamW} 则将权重衰减从梯度计算中分离出来——动量和二阶矩只基于数据损失的梯度计算，L2 正则化在最终更新时才加入。

\begin{knowledgebox}{为什么 AdamW 更受欢迎？}
在 Adam 中，L2 正则化的梯度会影响动量和自适应学习率的估计，这可能并非我们期望的行为。AdamW 让优化器专注于损失景观的导航，而正则化仅作为一个独立的权重缩减步骤。Meta 的 LLaMA 系列模型都使用 AdamW 优化器。
\end{knowledgebox}

\subsection{本章小结}

\begin{itemize}
    \item \textbf{SGD}：最基础的优化器，存在窄谷震荡、局部极小值和鞍点问题
    \item \textbf{SGD + Momentum}：通过速度累积解决上述问题
    \item \textbf{RMSProp}：通过逐元素缩放梯度，在平坦方向走大步、陡峭方向走小步
    \item \textbf{Adam}：结合 Momentum 和 RMSProp，加偏差修正，是最常用的优化器
    \item \textbf{AdamW}：解耦权重衰减，在许多实际任务中表现更好
\end{itemize}

%% ========== 学习率调度 ==========

